\chapter{Puissances}

\begin{itemize}
\it Calculer les puissances d'exposant $2$, $3$, $4$, $5$ des nombres suivants : 
\[ (+1); \phantom{miou}(+3); \phantom{miou}(+5); \phantom{miou}(-2); \phantom{miou}(-4); \phantom{miou}(-5).\]
\it Calculer : \[ (-11)^3; \phantom{miou}(+9)^2; \phantom{miou}(-10)^5; \phantom{miou}(+7)^4; \phantom{miou}(-13)^4.\]
\it Calculer : \[ \left(-\frac12\right)^5; \phantom{miou}\left(+\frac23\right)^2; \phantom{miou}\left(-\frac34\right)^3; \phantom{miou}(+1,4)^2; \phantom{miou}(-0,25)^2.\]
\it Calculer : \[ 5^3\times 5^2; \phantom{miou}(-7)\times(-7)^2; \phantom{miou}\left(\frac23\right)\times \left(\frac23\right)^3; \phantom{miou}
\left[\left(-\frac14\right)^2\right]^2\left[(-2)^2\right]^3\]
\item Calculer : 
\begin{itemize}
\it $[(-1)(+2)(-3)]^2$.
\it $[(+1)(-3)(-5)]^3$.
\it $[(+2)(-2)(-3)]^2$.
\it $\displaystyle\left[\left(\frac12\right)\left(-\frac13\right)\left(-\frac15\right)\right]^2$.
\it $ \displaystyle\left[(-5)\left(-\frac35\right)
\left(+\frac13\right)\right]^3$.
\it $\displaystyle \left[ (-0,5)(+0,75)\left(+\frac12\right)\right]^3$.
\end{itemize}
\it Calculer : 
\[ (-5)^7\div (-5)^3; \phantom{miaou}
(+0,01)^3\div (+0,01); \phantom{miaou} 
\left(-\frac45\right)^2\div \left(-\frac45\right)^2.\]
\[ 4^3\div 4^6; \phantom{miaou} (-2,5)^2\div (-2,5)^3; 
\phantom{miaou} \left(\frac23\right)^3\div \left(\frac23\right)^6.\]
\it Calculer la valeur de $x^4-x^3+x^2-x+1$. 
\begin{enumerate}
\item pour $x=3$;
\item pour $x=-2$;
\item pour $x=\frac13$;
\item pour $x=-\frac12$.
\end{enumerate}
\it Calculer la valeur de $x^3-3x^2+5x-7$. 
\begin{enumerate}
\item pour $x=4$;
\item pour $x=-1$;
\item pour $x=\frac23$;
\item pour $x=-\frac23$.
\end{enumerate}
\it Calculer la valeur de $\displaystyle\frac{x^3-1}{x^2+1}$. 
\begin{enumerate}
\item pour $x=1$;
\item pour $x=-1$;
\item pour $x=\frac14$;
\item pour $x=-\frac14$.
\end{enumerate}
\it Calculer la valeur de $a^4-4a^3b+6a^2b^2-4ab^3+b^4$. 
\begin{enumerate}
\item pour $a=2$ et $b=-3$;
\item pour $a=-\frac12$ et $b=-\frac12$.
\end{enumerate}
\it Calculer la valeur de $\displaystyle\frac{3xy}{x^2+y^2}-frac{x-y}{x+y}$.
\begin{enumerate}
\item pour $x=1$ et $y=-2$;
\item pour $x=-\frac12$ et $y=+\frac14$.
\end{enumerate}
\it Soient $3$ points $A$, $B$, $C$ d'abscisses $+2$, $-1$, $-3$, sur un axe où $O$ est l'origine des abscisses. Vérifier la relation : 
\[ \overline{OA}^2.\overline{BC}+\overline{OB}^2.\overline{CA}+\overline{OC}^2.\overline{AB} + \overline{AB}.\overline{BC}.\overline{CA} = 0.\]
\it Sur un axe où $O$ est l'origine des abscisses, sont placés trois points $A$, $B$, $M$ d'abscisses $+3$, $-3$, $+1$. Vérifier les relations : 
\[\overline{MA}.\overline{MB}=\overline{MO}^2-\overline{OA}^2.\]
\[ \overline{MA}^2+\overline{MB}^2=2\overline{MO}^2+2\overline{OA}^2.\]
\[\overline{MA}^2-\overline{MB}^2=2\overline{AB}.\overline{OM}.\]
\end{itemize}